\documentclass[UTF8,a4paper]{ctexart}
\usepackage{amsmath,amssymb}
\usepackage{graphicx}
\usepackage{array}
\usepackage{multirow}
\usepackage{geometry}
\usepackage{enumitem}

\usepackage{fancyhdr}
\usepackage{lastpage}


\usepackage{setspace}


\DeclareMathOperator{\arccot}{arccot}
\newcommand{\dlim}{\displaystyle\lim}
\newcommand{\dint}{\displaystyle\int}
\newcommand{\e}{\mathrm{e}}
\renewcommand{\i}{\mathrm{i}}


\usepackage{tikz}
\usetikzlibrary{calc}
\usepackage{eso-pic}


\geometry{left=31.7mm, right=31.7mm, top=30.4mm, bottom=30.4mm}



\pagestyle{fancy}
\fancyhead{}
\fancyfoot{}
\fancyfoot[C]{\normalsize 数学试题第\thepage 页（共 \pageref{LastPage} 页）}
\renewcommand{\headrulewidth}{0pt}

\begin{document}


\AddToShipoutPicture*{%
\begin{tikzpicture}[remember picture,overlay]
\draw[dashed,line width=0.4pt] ($(current page.north west)+(1.2cm,0)$) -- ($(current page.south west)+(1.2cm,0)$);
\end{tikzpicture}}



\AddToShipoutPicture*{%
\begin{tikzpicture}[remember picture,overlay]
% 信息栏放在装订线左侧（更靠左），装订后会被遮挡
% 在页面左侧垂直居中，竖排文字：姓名_学号_班级
% 页面高度约29.7cm，居中位置约14.8cm
% 从下至上排列：姓名 → 学号 → 班级
% 每个项目占3cm高度
% 姓名（最下方）
\node[anchor=west,font=\small\bfseries,rotate=90] at ($(current page.south west)+(0.5cm,6cm)$) {姓名\underline{\hspace{2.5cm}}};
% 学号（中间）
\node[anchor=west,font=\small\bfseries,rotate=90] at ($(current page.south west)+(0.5cm,10cm)$) {学号\underline{\hspace{2.5cm}}};
% 班级（最上方，可选）

\node[anchor=west,font=\small\bfseries,rotate=90] at ($(current page.south west)+(0.5cm,14cm)$) {班级\underline{\hspace{2cm}}};

\end{tikzpicture}}


\setstretch{1.5}


\fontsize{11}{13.2}\selectfont




\begin{center}

{\textbf{\fontsize{22pt}{26.4pt}\selectfont 期中考试}}
\end{center}

\vspace{1em}

\noindent\textbf{注意事项：}\par\smallskip
1. 答卷前，考生务必将自己的姓名、准考证号填写在答题卡上。\par
2. 回答选择题时，选出每小题答案后，用铅笔把答题卡上对应题目的答案标号涂黑。如需改动，用橡皮擦干净后，再选涂其他答案标号。回答非选择题时，将答案写在答题卡上。写在本试卷上无效。\par
3. 考试结束后，将本试卷和答题卡一并交回。

\par\vspace{6pt}



\noindent\textbf{一、选择题：本题共8小题，每小题5分，共40分。在每小题给出的四个选项中，只有一项是符合题目要求的。}\hangindent=2em
\par\medskip




\begin{enumerate}[series=s0, leftmargin=2em, itemsep=0.8em]

\end{enumerate}
\vspace{1em}



\noindent\textbf{二、选择题：本题共3小题，每小题6分，共18分。在每小题给出的选项中，有多项符合题目要求。全部选对的得6分，部分选对的得部分分，有选错的得0分。}\hangindent=2em
\par\medskip




\begin{enumerate}[series=s1, leftmargin=2em, itemsep=0.8em]

\end{enumerate}
\vspace{1em}



\noindent\textbf{三、填空题：本题共3小题，每小题5分，共15分。}\hangindent=2em
\par\medskip




\begin{enumerate}[series=s2, leftmargin=2em, itemsep=0.8em]

\end{enumerate}
\vspace{1em}



\noindent\textbf{四、解答题：本题共5小题，共77分。解答应写出文字说明、证明过程或演算步骤。}\hangindent=2em
\par\medskip




\begin{enumerate}[series=s3, leftmargin=2em, itemsep=0.8em]

\end{enumerate}
\vspace{1em}


\end{document}